\scanpage{149}\mspage{23}
Venons-en à la catégorie $M^r$ elle $=$. L'argument de la page 21 montre que l'ens des \inserted{classes d'isom} objets de $M^r$ indexés par un $I$ fixé, de cardinal $\le r^{\pi_0}$, comme
\[
\deleted{\unsure{X|A}}A=\underset{i}{\underset{\unsure{\mathrel{\text{---}}|\longrightarrow}}{\lim}}A_i
\quad\text{avec }A_i\in\mathrm{Filt}^{\pi_0},
\]
est de cardinal majoré par $r^{\pi_0}$ (si $r\ge C_0$). Si on fait varier $I$, on trouve \deleted{\unsure{la|le}} une majoration par $r^{\pi_0}\cdot2^{(r^{\pi_0})}=2^{(r^{\pi_0})}$.
On en \unsure{conclut} comme \unsure{tantôt|tantot} que \unsure{l'ens|bien} $\widetilde{\mathrm{Fl}\,M^r}$ est de cardinal majoré par
\[
\left(2^{(r^{\pi_0})}\right)^2\times r^{(r^{\pi_0})}\qquad(\text{si }r\ge\mathrm{Sup}(\pi_1,C_1)),
\]
donc par $r^{(r^{\pi_0})}$. Je dis qu'on a
\[
r^{(r^{\pi_0})}=2^{(r^{\pi_0})},
\]
en effet, on a $\ge$, et pour prouver $\le$, on note que $r<2^r$, donc
\[
r^{(r^{\pi_0})}\le(2^r)^{(\unsure{r}^{\pi_0})}=2^{r\cdot r^{\pi_0}}=2^{r^{\pi_0}},
\]
\unsure{\deleted{Cet argument}|Cet argument}
\deleted{\unsure{On a établi}}
\unsure{OK|ok} ! Cela établit donc 3.7.
\unsure{On a} en fait le
\uline{Lemme} \uline{3.2.} \unsure{Soient|Soit} $r$ un cardinal, $r'$ un cardinal $\ge r$. Alors on a
\[
r^{r'}=2^{r'}\qquad(=s^{r'}\text{ pour \unsure{tous|tout|ts} }2\le s\le r)
\]
\uline{Corollaire 1} \unsure{Soient|Soit} \deleted{$r\ge s\ge\mathrm{Sup}(\pi_1,C_1)$.} $r\ge\mathrm{Sup}(\pi_1,C_1)$, et $A,X\in\mathrm{Filt}^{\unsure{r|\pi}}$. Alors
\[
\mathrm{card}\,\mathrm{Hom}_M(A,X)\le2^{(r^{\pi_0})}
\qquad(=2^r\text{ si }r=2^c,\ \inserted{\text{avec}}\ c\unsure{\ge|>}\pi_0\unsure{)}
\tag*{(3.10),}
\]
On applique 3.8., \deleted{\unsure{car|au}} qui donne la majoration $r^{(r^{\pi_0})}$, et le lemme 3.2.

\deleted{\uline{Corollaire 2}}
On \unsure{note|voit} qu'on trouve la même majoration que pour le cardinal de $\widetilde{\underline{\mathrm{Fl}}(M^r)}$ \unsure{démonstration|détermination} (3.7)! $\unsure{D|\Delta}\unsure{=|\simeq}$ (3.9.) donne

\uline{Corollaire 2} Soit $r\ge\mathrm{Sup}(\pi_1,C_1)$, et $\unsure{\alpha|X}\in\mathrm{Ob}\,M^r$. Alors
\[
\mathrm{card}\,\widetilde{\underline{\mathrm{Fl}}(M^r_{/X})}\le2^{(\unsure{r}^{\pi_0})}
\qquad(=2^r\text{ si }r=2^c,\text{ avec }c\ge\pi_0).
\tag*{(3.11).}
\]
c'est donc la formule (3.9) pour $s\le r$, \unsure{laissée|laissé|laissee} en suspens.
% archive sheet number: 148
